\documentclass[14pt]{extarticle}

%% Поля: левое 2 см, правое 1 см, верхнее 2 см, нижнее 2 см (требование ВШЭ)
\usepackage[left=2cm, right=1cm, top=2cm, bottom=2cm]{geometry}

%% Кодировка и язык
\usepackage[utf8]{inputenc}
\usepackage[T2A]{fontenc}
\usepackage[russian, english]{babel}

%% Шрифт Times New Roman (tempora — кириллический клон Times)
\usepackage{tempora}

%% Межстрочный интервал 1.5 (требование ВШЭ)
\usepackage{setspace}
\onehalfspacing

%% Абзацный отступ 1.25 см (требование ВШЭ)
\usepackage{indentfirst}
\setlength{\parindent}{1.25cm}

%% Нумерация страниц — внизу по центру (требование ВШЭ)
\usepackage{fancyhdr}
\pagestyle{fancy}
\fancyhf{}
\fancyfoot[C]{\thepage}
\renewcommand{\headrulewidth}{0pt}
\fancypagestyle{plain}{
    \fancyhf{}
    \fancyfoot[C]{\thepage}
    \renewcommand{\headrulewidth}{0pt}
}

%% Математика
\usepackage{amsmath}
\usepackage{amssymb}
\usepackage{bm}
\usepackage{siunitx}
\usepackage{cases}

%% Графика и таблицы
\usepackage{graphicx}
\usepackage{float}
\usepackage{caption}
\usepackage{subcaption}
\renewcommand\thesubfigure{\asbuk{subfigure}}
\usepackage{booktabs}
\usepackage{multirow}
\usepackage{array}
\usepackage{dcolumn}

%% Алгоритмы
\usepackage[ruled,vlined,linesnumbered]{algorithm2e}
\SetAlgorithmName{Алгоритм}{алгоритм}{Список алгоритмов}

%% Код Python
\usepackage{listings}
\usepackage{xcolor}

\definecolor{codebg}{HTML}{F7F7F7}
\definecolor{codegreen}{HTML}{2E7D32}
\definecolor{codeblue}{HTML}{1565C0}
\definecolor{codegray}{HTML}{757575}

\lstdefinestyle{python}{
    language=Python,
    backgroundcolor=\color{codebg},
    basicstyle=\ttfamily\small,
    keywordstyle=\color{codeblue}\bfseries,
    stringstyle=\color{codegreen},
    commentstyle=\color{codegray}\itshape,
    numberstyle=\tiny\color{codegray},
    breaklines=true,
    captionpos=b,
    numbers=left,
    numbersep=8pt,
    showstringspaces=false,
    tabsize=4,
    frame=single,
    framerule=0.4pt,
    rulecolor=\color{codegray!40},
    xleftmargin=15pt,
}
\lstset{style=python}

%% Прочее оформление
\usepackage{enumitem}
\usepackage{hyphenat}
\hyphenation{ма-те-ма-ти-ка вос-ста-нав-ли-вать се-ман-ти-чес-кая сег-мен-та-ция}
\setlength{\emergencystretch}{30pt}

%% Гиперссылки
\usepackage{hyperref}
\definecolor{linkcolor}{HTML}{0000D5}
\hypersetup{
    pdfstartview=FitH,
    linkcolor=linkcolor,
    urlcolor=linkcolor,
    citecolor=linkcolor,
    colorlinks=true,
    pdftitle={Оценка риска повреждения фасадов зданий},
    pdfauthor={Фамилия Имя Отчество},
}

%% Библиография
\bibliographystyle{IEEEtran}

\newcommand{\miou}{mIoU\xspace}

% ---------------------------------------------------------------
\begin{document}

%% ===== ТИТУЛЬНЫЙ ЛИСТ (без нумерации) =====
\begin{titlepage}
    \newgeometry{left=1.5cm, right=1.5cm, top=2cm, bottom=2cm}
    \thispagestyle{empty}
    \begin{center}
        {\normalsize ФЕДЕРАЛЬНОЕ ГОСУДАРСТВЕННОЕ АВТОНОМНОЕ ОБРАЗОВАТЕЛЬНОЕ УЧРЕЖДЕНИЕ\\
        ВЫСШЕГО ОБРАЗОВАНИЯ}\\
        \vspace{0.2cm}
        {\normalsize «НАЦИОНАЛЬНЫЙ ИССЛЕДОВАТЕЛЬСКИЙ УНИВЕРСИТЕТ\\
        «ВЫСШАЯ ШКОЛА ЭКОНОМИКИ»»}\\
        \vspace{1cm}
        {\normalsize Факультет компьютерных наук}\\
        \vspace{1.5cm}

        {\normalsize Фамилия Имя Отчество}\\
        \vspace{0.2cm}
        {\normalsize\textbf{Оценка риска повреждения фасадов городских зданий\\
        на основе семантической сегментации и графовых нейронных сетей}}\\
        \vspace{0.2cm}
        {\normalsize Выпускная квалификационная работа}\\
        {\normalsize по направлению подготовки 09.04.04 «Программная инженерия»}\\
        {\normalsize образовательная программа «Искусственный интеллект»}\\
    \end{center}

    \vspace{4cm}
    \hfill
    \begin{minipage}{0.45\textwidth}
        \begin{flushleft}
            {\normalsize Научный руководитель}\\
            {\normalsize должность, учёная степень}\\
            {\normalsize Фамилия И.\,О.}\\
        \end{flushleft}
    \end{minipage}

    \vspace{4cm}
    \begin{center}
        Санкт-Петербург 2025
    \end{center}
    \restoregeometry
\end{titlepage}

%% ===== АННОТАЦИЯ (800–2000 знаков) =====
\section*{Аннотация}
\addcontentsline{toc}{section}{Аннотация}

Данная работа посвящена разработке автоматизированного конвейера оценки риска повреждения
фасадов городских зданий по данным уличных панорам. Объектом исследования являются
фасады жилых зданий Санкт-Петербурга. Предлагаемый подход объединяет три этапа:
(1)~семантическую сегментацию дефектов на основе трансформерной архитектуры SegFormer-B2,
(2)~ранжирование зданий по степени повреждённости с помощью модели попарных предпочтений
RankNet, (3)~построение графа зданий и обучение графовой нейронной сети (GNN) для итоговой
оценки риска. Набор данных сформирован из панорамных снимков Яндекс.Панорам и включает
три суперкласса дефектов: антропогенные (граффити), биохимические (плесень, высолы) и
механические (трещины). Дополнительно используется набор профессиональных DSLR-фотографий
для оценки межсъёмочной обобщаемости модели. Базовая модель сегментации достигает
mIoU~=~0,595. Работа содержит [NN]~страниц, [NN]~рисунков, [NN]~таблиц,
[NN]~использованных источников.

\vspace{0.5em}
\noindent\textbf{Ключевые слова:} семантическая сегментация, графовые нейронные сети,
оценка технического состояния зданий, SegFormer, RankNet, компьютерное зрение,
Яндекс.Панорамы.

\newpage

%% ===== ANNOTATION =====
\section*{Annotation}
\addcontentsline{toc}{section}{Annotation}

This work presents an automated pipeline for assessing facade damage risk in urban buildings
using street-level panoramic imagery. The proposed approach combines three stages:
(1)~semantic segmentation of defects using the SegFormer-B2 transformer architecture,
(2)~building-level damage ranking via the RankNet pairwise preference model, and
(3)~a building graph construction with a Graph Neural Network (GNN) for final risk scoring.
The dataset is collected from Yandex Panoramas and annotated with three defect superclasses:
anthropological (graffiti), biochemical (mold, efflorescence), and mechanical (cracks).
An additional set of professional DSLR photographs is used to evaluate cross-domain
generalization. The baseline segmentation model achieves mIoU~=~0.595. The work contains
[NN]~pages, [NN]~figures, [NN]~tables, [NN]~references.

\vspace{0.5em}
\noindent\textbf{Keywords:} semantic segmentation, graph neural networks, building condition
assessment, SegFormer, RankNet, computer vision, Yandex Panoramas.

\newpage

%% ===== СОДЕРЖАНИЕ =====
\tableofcontents
\newpage

% ---------------------------------------------------------------
%% ===== ВВЕДЕНИЕ (не более 10% объёма ВКР) =====
% ---------------------------------------------------------------
\section*{Введение}
\addcontentsline{toc}{section}{Введение}
\label{sec:intro}

Мониторинг технического состояния зданий является важной задачей городского управления.
Регулярные обследования фасадов позволяют своевременно выявлять дефекты и предотвращать
аварийные ситуации. Ручной осмотр трудоёмок, дорогостоящ и не поддаётся масштабированию
на уровень города.

\textbf{Актуальность.} ...

\textbf{Научная новизна.} ...

\textbf{Цель работы} --- разработать автоматизированный конвейер оценки риска повреждения
фасадов городских зданий по данным уличных панорам.

\noindent\textbf{Задачи:}
\begin{enumerate}[leftmargin=*, label=\arabic*.]
    \item Сформировать набор данных панорамных снимков фасадов зданий с разметкой
          трёх классов дефектов.
    \item Обучить модель семантической сегментации дефектов (SegFormer-B2).
    \item Построить датасет попарных сравнений зданий и обучить модель ранжирования (RankNet).
    \item Построить граф зданий и обучить графовую нейронную сеть (GNN) для оценки риска.
    \item Оценить качество каждого этапа конвейера и провести анализ ошибок.
\end{enumerate}

\noindent\textbf{Объект исследования} --- фасады жилых зданий Санкт-Петербурга.

\noindent\textbf{Методы исследования} --- глубокое обучение, семантическая сегментация,
обучение ранжированию, графовые нейронные сети.

В главе~1 приводится обзор смежных работ и формулируется постановка задачи.
Глава~2 описывает набор данных, предложенную архитектуру и методы обучения.
В главе~3 представлены результаты экспериментов и рекомендации по применению.

\newpage

% ---------------------------------------------------------------
%% ===== ГЛАВА 1. Обзор литературы и постановка задачи =====
% ---------------------------------------------------------------
\section{Обзор литературы и постановка задачи}
\label{sec:review}

\subsection{Автоматизированная диагностика дефектов зданий}

\subsection{Семантическая сегментация в задачах строительного контроля}

\subsection{Ранжирование объектов в задачах оценки состояния}

\subsection{Графовые нейронные сети для городских данных}

\subsection{Постановка задачи}

\subsection{Выводы по главе 1}
% Требование ВШЭ: 2–5 пунктов
\begin{enumerate}[leftmargin=*, label=\arabic*.]
    \item ...
    \item ...
    \item ...
\end{enumerate}

\newpage

% ---------------------------------------------------------------
%% ===== ГЛАВА 2. Данные и предлагаемый метод =====
% ---------------------------------------------------------------
\section{Данные и предлагаемый метод}
\label{sec:method}

\subsection{Набор данных}

\subsubsection{Источники данных}

Набор данных сформирован из двух источников:
\begin{enumerate}[leftmargin=*, label=\arabic*.]
    \item \textbf{Яндекс.Панорамы} --- скриншоты фасадов $\approx71$~здания,
          нарезанные на патчи $800{\times}800$ пикселей ($\approx206$~патчей).
    \item \textbf{DSLR-фотографии} --- $\approx500$~профессиональных снимков зданий
          в районе станции метро «Технологический институт», используемые как
          тест-сет для оценки доменной обобщаемости.
\end{enumerate}

\subsubsection{Классы дефектов}

В работе выделяются три суперкласса дефектов:
\begin{itemize}[leftmargin=*]
    \item \textbf{Антропогенные} (Anthropological, A) --- граффити, следы вандализма.
    \item \textbf{Биохимические} (Bio-chemical, B) --- плесень, высолы, биологические
          отложения.
    \item \textbf{Механические} (Mechanical, M) --- трещины, сколы, разрушение материала.
\end{itemize}

\subsubsection{Разметка}

Разметка выполнена в системе CVAT с применением автоматической предразметки
SAM2~\cite{kirillov2023sam}.

\subsection{Семантическая сегментация (SegFormer-B2)}

SegFormer~\cite{xie2021segformer} --- трансформерная архитектура для семантической
сегментации, состоящая из иерархического энкодера Mix Transformer (MiT-B2) и лёгкого
MLP-декодера.

Для борьбы с дисбалансом классов применяется взвешенная кросс-энтропия:
\begin{equation}
    \mathcal{L} = -\sum_{c} w_c \sum_{i} y_{ic} \log \hat{p}_{ic},
    \label{eq:loss}
\end{equation}
где $w_c$ --- вес класса $c$, $y_{ic}$ --- истинная метка пикселя $i$,
$\hat{p}_{ic}$ --- предсказанная вероятность.

\subsection{Ранжирование зданий (RankNet)}

\subsection{Граф зданий и GNN}

\subsection{Выводы по главе 2}
% Требование ВШЭ: 3–5 пунктов, элементы научной новизны
\begin{enumerate}[leftmargin=*, label=\arabic*.]
    \item ...
    \item ...
    \item ...
\end{enumerate}

\newpage

% ---------------------------------------------------------------
%% ===== ГЛАВА 3. Результаты и рекомендации по применению =====
% ---------------------------------------------------------------
\section{Результаты и рекомендации по применению}
\label{sec:results}

\subsection{Условия проведения экспериментов}

\subsection{Метрики качества}

Основной метрикой сегментации является mean Intersection over Union (\miou):
\begin{equation}
    \text{mIoU} = \frac{1}{C} \sum_{c=1}^{C}
    \frac{TP_c}{TP_c + FP_c + FN_c},
    \label{eq:miou}
\end{equation}
где $C$ --- число классов, $TP_c$, $FP_c$, $FN_c$ --- истинно-положительные,
ложно-положительные и ложно-отрицательные предсказания для класса $c$.

\subsection{Результаты сегментации}

Базовая модель SegFormer-B2 достигает \miou~$= 0{,}595$ на валидационной
выборке (табл.~\ref{tab:seg_results}).

\begin{table}[H]
    \centering
    \caption{Результаты семантической сегментации по классам.}
    \label{tab:seg_results}
    \begin{tabular}{lccc}
        \toprule
        \textbf{Класс} & \textbf{IoU} & \textbf{Precision} & \textbf{Recall} \\
        \midrule
        Антропогенные  & 0{,}66 & --- & 0{,}98 \\
        Биохимические  & 0{,}45 & --- & ---    \\
        Механические   & 0{,}46 & --- & 0{,}57 \\
        \midrule
        \textbf{mIoU}  & \textbf{0{,}595} & & \\
        \bottomrule
    \end{tabular}
\end{table}

\subsection{Анализ ошибок}

\subsection{Результаты ранжирования}

\subsection{Результаты GNN}

\subsection{Рекомендации по применению}

\newpage

% ---------------------------------------------------------------
%% ===== ЗАКЛЮЧЕНИЕ (5–7 пунктов) =====
% ---------------------------------------------------------------
\section*{Заключение}
\addcontentsline{toc}{section}{Заключение}

\begin{enumerate}[leftmargin=*, label=\arabic*.]
    \item ...
    \item ...
    \item ...
    \item ...
    \item ...
\end{enumerate}

\newpage

% ---------------------------------------------------------------
%% Перечень сокращений (опционально, раскомментировать при необходимости)
% ---------------------------------------------------------------
%\section*{Перечень сокращений}
%\addcontentsline{toc}{section}{Перечень сокращений}
%\begin{tabular}{ll}
%    GNN       & Graph Neural Network (графовая нейронная сеть) \\
%    mIoU      & mean Intersection over Union                   \\
%    SegFormer & Segmentation Transformer                       \\
%    RankNet   & Ranking Neural Network                         \\
%\end{tabular}
%\newpage

% ---------------------------------------------------------------
%% ===== СПИСОК ИСПОЛЬЗОВАННЫХ ИСТОЧНИКОВ =====
% ---------------------------------------------------------------
\addcontentsline{toc}{section}{Список использованных источников}
\bibliography{bib}

\end{document}
